\section{Desenvolvimento}
\label{sec:desenvolvimento}

Esta e a parte principal do artigo. Divida-a em quantas secoes forem
necessarias: referencial teorico, metodologia, resultados e discussao. Cada
afirmacao que nao seja sua precisa vir acompanhada da fonte, como em
\cite{exemplo2024}, e obras de referencia da area entram do mesmo modo
\cite{knuth1984}.

\subsection{Figuras}
\label{subsec:figuras}

Toda figura precisa de legenda acima e indicacao de fonte abaixo, conforme a
NBR 14724. Coloque os arquivos de imagem em \texttt{figures/}.

% Exemplo de figura. Descomente depois de colocar a imagem em figures/.
%
% \begin{figure}[htb]
%   \centering
%   \caption{Descricao breve do que a figura mostra}
%   \label{fig:exemplo}
%   \includegraphics[width=0.7\textwidth]{exemplo.png}
%   \fonte{Elaborado pelo autor (2026).}
% \end{figure}

\subsection{Tabelas}
\label{subsec:tabelas}

Tabelas seguem o padrao do IBGE: sem linhas verticais e abertas nas laterais.

\begin{table}[htb]
  \centering
  \caption{Exemplo de tabela no padrao do IBGE}
  \label{tab:exemplo}
  \begin{tabular}{lrr}
    \toprule
    \textbf{Categoria} & \textbf{Frequencia} & \textbf{Percentual} \\
    \midrule
    Primeira categoria & 120 & 40,0\% \\
    Segunda categoria  &  90 & 30,0\% \\
    Terceira categoria &  90 & 30,0\% \\
    \midrule
    Total              & 300 & 100,0\% \\
    \bottomrule
  \end{tabular}
  \fonte{Elaborado pelo autor (2026).}
\end{table}

\subsection{Equacoes}
\label{subsec:equacoes}

Equacoes deslocadas do texto sao numeradas e podem ser referenciadas, como a
Equacao~\eqref{eq:exemplo}:

\begin{equation}
  \label{eq:exemplo}
  \bar{x} = \frac{1}{n} \sum_{i=1}^{n} x_i
\end{equation}
